\origpage{192}{187}
\begin{statement}{引理 11.3}
固定 $k$ 到 $\mathbb C$ 的一个嵌入。则\term{除子类}{divisor class}{classe de diviseurs}
\[
\mathscr P_{ij}:=(\operatorname{pr}_i+\operatorname{pr}_j)^*L-\operatorname{pr}_i^*L-\operatorname{pr}_j^*L
\]
所对应的 $H^{1,1}_{\bar\partial}(A^n)$ 中的\term{上同调类}{cohomology class}{classe de cohomologie}，在 $A(\mathbb C)^n$ 上可由下列空间中的一个\term{微分形式}{differential form}{forme différentielle}表示：
\[
\operatorname{pr}_i^*\mathscr E^{1,0}(A)\otimes\operatorname{pr}_j^*\mathscr E^{0,1}(A)
+\operatorname{pr}_i^*\mathscr E^{0,1}(A)\otimes\operatorname{pr}_j^*\mathscr E^{1,0}(A)
\subseteq\mathscr E^{1,1}(A^n).
\]
\end{statement}
\begin{proof}
设 $d=\dim A$，利用 $A$ 是 $\mathbb C^d$ 模去一个格所得的商这一表示，在 $A$ 上取局部坐标 $z_1,\ldots,z_d$。给 $\mathcal O(L)$ 赋予一个曲率平移不变的度量，其曲率可写成
\[
\sum_{\alpha,\beta=1}^d a_{\alpha\beta}\,dz_\alpha\wedge d\bar z_\beta,
\]
其中 $a_{\alpha\beta}$ 是常数。对 $\alpha=1,\ldots,d$，令 $u_\alpha=\operatorname{pr}_i^*z_\alpha$，$v_\alpha=\operatorname{pr}_j^*z_\alpha$。使用上述 $L$ 上的度量，$\mathscr P_{ij}$ 的曲率就是
\[
\sum_{\alpha,\beta=1}^d a_{\alpha\beta}
(du_\alpha\wedge d\bar v_\beta+dv_\alpha\wedge d\bar u_\beta).
\]
\end{proof}
比较次数，便立即得到下面的主要齐次性引理。
\begin{statement}{推论 11.4}
在 $\prod X_i$ 上，任一具有最大\term{余维}{codimension}{codimension}的\term{交积}{intersection product}{produit d'intersection}
\[
\prod_{i<j}\mathscr P_{ij}^{e_{ij}}\cdot\prod_{i=1}^n\operatorname{pr}_i^*L^{e_i}
\]
均为零，除非
\[
2e_i+\sum_{j<i}e_{ji}+\sum_{j>i}e_{ij}=2\dim X_i,
\qquad i=1,\ldots,n.
\]
因此，由于
\[
(s_i\cdot\operatorname{pr}_i-s_j\cdot\operatorname{pr}_j)^*L
=s_i^2\operatorname{pr}_i^*L+s_j^2\operatorname{pr}_j^*L+s_is_j\mathscr P_{ij},
\]
$L_{\delta,s}$ 在 $\prod X_i$ 上的最高次\term{自交数}{self-intersection number}{nombre d'auto-intersection}对每个 $s_i$ 都是 $2\dim X_i$ 次齐次的。
\end{statement}
下一命题是本节的主要结果，其证明由 M.~Nakamaye 和 G.~Faltings 分别独立给出。
\begin{statement}{命题 11.5}
存在仅依赖于 $X,A,L,\dim X_1,\ldots,\dim X_n$ 以及 $\deg X_i$ 的上界的常数 $c>0$ 和 $\epsilon>0$，使得对每个正有理数元组 $s=(s_1,\ldots,s_n)$，当 $d$ 充分大（可依赖于 $s$）时，均有
\[
h^0\!\left(\prod X_i,dL_{-\epsilon,s}\right)
>c d^{\sum\dim X_i}\prod_{i=1}^n s_i^{2\dim X_i}.
\]
\end{statement}
\origpage{193}{188}
\begin{proof}
由 Seshadri 判据（[H 1]，第一章，§7），$\delta>0$ 时 $L_{\delta,s}$ 是\term{丰沛}{ample}{ample}的。由 Riemann--Roch 定理，
\[
h^0\!\left(\prod X_i,dL_{\delta,s}\right)
=d^{\sum\dim X_i}\frac{(L_{\delta,s}^{\sum\dim X_i})}{(\sum\dim X_i)!}(1+o(1)).
\]
为简化记号，在本证明余下部分令 $N=\sum\dim X_i$。对每个指标 $j$，令 $H_j$ 是由 $\operatorname{pr}_j^*\mathcal O(L)$ 的某个截面在 $\prod X_i$ 上截出的子概形。与上面一样，
\[
h^0(H_j,dL_{\delta,s}|_{H_j})
=d^{N-1}\frac{(\operatorname{pr}_j^*L\cdot L_{\delta,s}^{N-1})}{(N-1)!}(1+o(1)).
\]
对每个 $i$（包括 $j$），$\operatorname{pr}_i^*\mathcal O(L)$ 在 $H_j$ 上都可由一个\term{有效除子}{effective divisor}{diviseur effectif}表示，故对每个非负整数元组 $m_1,\ldots,m_n$，
\[
h^0\!\left(H_j,\left.(dL_{\delta,s}-\sum m_i\operatorname{pr}_i^*L)\right|_{H_j}\right)
\le d^{N-1}\frac{(\operatorname{pr}_j^*L\cdot L_{\delta,s}^{N-1})}{(N-1)!}(1+o(1)).
\]
注意，这里的 $o(1)$ 不依赖于 $m_1,\ldots,m_n$。\term{短正合列}{short exact sequence}{suite exacte courte}
\begin{align*}
0&\longrightarrow\mathcal O\!\left(dL_{\delta,s}-\sum m_i\operatorname{pr}_i^*L-\operatorname{pr}_j^*L\right)
\longrightarrow\mathcal O\!\left(dL_{\delta,s}-\sum m_i\operatorname{pr}_i^*L\right)\\
&\longrightarrow\left.\mathcal O\!\left(dL_{\delta,s}-\sum m_i\operatorname{pr}_i^*L\right)\right|_{H_j}\longrightarrow0
\end{align*}
所诱导的上同调正合列给出
\begin{align*}
h^0\!\left(\prod X_i,dL_{\delta,s}-\sum m_i\operatorname{pr}_i^*L-\operatorname{pr}_j^*L\right)
&\ge h^0\!\left(\prod X_i,dL_{\delta,s}-\sum m_i\operatorname{pr}_i^*L\right)\\
&\quad-h^0\!\left(H_j,dL_{\delta,s}-\sum m_i\operatorname{pr}_i^*L\right)\\
&\ge h^0\!\left(\prod X_i,dL_{\delta,s}-\sum m_i\operatorname{pr}_i^*L\right)\\
&\quad-d^{N-1}\frac{(\operatorname{pr}_j^*L\cdot L_{\delta,s}^{N-1})}{(N-1)!}(1+o(1)).
\end{align*}
因为 $dL_{-\epsilon,s}=dL_{\delta,s}-d(\delta+\epsilon)\sum_i s_i^2H_i$，上述估计遂给出
\begin{align*}
h^0\!\left(\prod X_i,dL_{-\epsilon,s}\right)
&\ge h^0\!\left(\prod X_i,dL_{\delta,s}\right)
-d(\delta+\epsilon)\sum_{i=1}^n s_i^2 h^0(H_i,dL_{\delta,s}|_{H_i})\\
&\ge d^N\left(\frac{(L_{\delta,s}^N)}{N!}
-(\delta+\epsilon)\sum_{i=1}^n s_i^2\frac{(\operatorname{pr}_i^*L\cdot L_{\delta,s}^{N-1})}{(N-1)!}\right)(1+o(1)).
\end{align*}
由推论 11.4，这个下界等于
\[
d^N\prod_{i=1}^n s_i^{2\dim X_i}
\left(\frac{(L_{\delta,1}^N)}{N!}
-(\delta+\epsilon)\sum_{i=1}^n\frac{(\operatorname{pr}_i^*L\cdot L_{\delta,1}^{N-1})}{(N-1)!}\right)(1+o(1)).
\]
\origpage{194}{189}
括号内是关于 $\delta$ 和 $\epsilon$ 的多项式；由推论 11.2，其常数项为正。因此可取充分小的 $\delta>0$、$\epsilon>0$ 和 $c>0$，使得
\[
0<c<\frac{(L_{\delta,1}^N)}{N!}
-(\delta+\epsilon)\sum_{i=1}^n\frac{(\operatorname{pr}_i^*L\cdot L_{\delta,1}^{N-1})}{(N-1)!}.
\]
这些交数是在变化的 $X_i$ 上取的，但这并不会造成实质困难：由于 $X_i$ 的次数有界，它们只能属于有限多个\term{数值等价类}{numerical equivalence classes}{classes d'équivalence numérique}。
\end{proof}

\section*{§12. 数的几何的进一步结果}
\addcontentsline{toc}{section}{§12. 数的几何的进一步结果}
设 $\Gamma$ 是\term{数域}{number field}{corps de nombres} $k$ 的\term{整数环}{ring of integers}{anneau des entiers} $R$ 上一个带度量的有限生成格。对 $k$ 的所有\term{阿基米德赋值}{archimedean place}{place archimédienne} $v$，在 $\Gamma$ 关于 $v$ 的\term{完备化}{completion}{complété} $\Gamma_v$ 上给定测度，使下述条件成立：若 $\gamma_1,\ldots,\gamma_\delta$ 是 $\Gamma$ 中 $R$ \term{线性无关}{linearly independent}{linéairement indépendant}且模去挠元后生成 $\Gamma$ 的一组向量，则它们在各 $v$ 处生成的 $\mathbb Z$ 格（$k_v\cong\mathbb R$ 时）或 $\mathbb Z[i]$ 格（$k_v\cong\mathbb C$ 时）的\term{余体积}{covolumes}{covolumes}之积，等于 $(\Gamma:\sum_{i=1}^\delta R\gamma_i)$。又令 $V(\Gamma)$ 为所有 $v$ 处 $\Gamma_v$ 中单位球（相对于 $\|\cdot\|_v$）的体积（相对于上述测度）之积。定义长度函数 $\ell(\gamma)=\prod_v\|\gamma\|_v$，并对 $i=1,\ldots,\delta$ 定义\term{逐次极小值}{successive minima}{minima successifs} $\lambda_i$ 为满足下述条件的最小 $\lambda$：存在 $R$ \term{线性无关}{linearly independent}{linéairement indépendant}的元素 $\gamma_1,\ldots,\gamma_i\in\Gamma$，且对 $j=1,\ldots,i$ 均有 $\ell(\gamma_j)\le\lambda$。
\begin{statement}{引理 12.1}
在上述情形下：
\begin{enumerate}[label=(\alph*).]
\item 存在仅依赖于 $k$ 的常数 $c_1,c_2$，使得
\[
(c_2\delta)^{-\delta}\le\lambda_1\cdots\lambda_\delta V(\Gamma)\le c_3^\delta.
\]
\item 设 $\beta:\Gamma_1\to\Gamma_2$ 是带度量的 $R$ 模同态。分别以 $\delta_0$、$\delta_2$ 表示 $\beta$ 的核与像作为 $R$ 模的秩。再设常数 $C$ 满足
\begin{equation*}\tag{12.2}
\prod_{v\mid\infty}\|\beta(\gamma)\|_v\le C\prod_{v\mid\infty}\|\gamma\|_v,
\qquad\gamma\in\Gamma_1.
\end{equation*}
则
\begin{equation*}\tag{12.3}
V(\Gamma_1)\le2^{[k:\mathbb Q]\delta_0}V(\Ker\beta)C^{\delta_2}V(\im\beta).
\end{equation*}
\end{enumerate}
\end{statement}
\begin{proof}
(a) 是 Minkowski 凸体理论的推广，可由（[V 1]，6.1.11）推出。该处仅对平行多面体证明了这一结论，但证明同样适用于更一般的凸长度函数。(b) 的证明是直接的。
\end{proof}
注意，此引理可以用来证明\term{数域}{number field}{corps de nombres}上的 Siegel 引理，虽然所得结果没有定理 2.4 那样精确。具体而言，在 $\Gamma_1:=R^N$ 上赋予度量：在实赋值处令 $\|x\|^2=|x_1|^2+\cdots+|x_N|^2$，在复赋值处令 $\|x\|=|x_1|^2+\cdots+|x_N|^2$。用定理 2.4 的语言，$h(x_i)=(1/[k:\mathbb Q])\log\lambda_i$。又在 (12.2) 中取
\[
C=\prod_{v\mid\infty}N^{[k_v:\mathbb R]}\max_{i,j}\|a_{ij}\|_v;
\]
它起到引理 2.1 中 $Q$ 的作用，或定理 2.4 中 $\exp(h(A))$ 的作用。
\origpage{195}{190}
这时 $V(\Gamma_1)$ 和 $V(\im\beta)$ 都只取决于相应的秩，因此 (12.3) 给出了仅依赖于 $C$ 和各个秩的 $V(\Ker\beta)$ 的下界。再由引理 (a)，得到 $\sum_{i=1}^{N-M}h(x_i)$ 的上界。

我们将利用引理 12.1 确定 $V(\Ker\beta)$ 的一个下界。只要能给出 $\delta_0$ 的合适下界，引理 (a) 就能推出存在一个\term{整体截面}{global section}{section globale} $\gamma$，它在 $\Gamma_1$ 中的度量很小（以 $\exp(cd\sum s_i^2)$ 为界）。这将在 §14 完成。

\section*{§13. Faltings 复形的算术}
\addcontentsline{toc}{section}{§13. Faltings 复形的算术}
为取得 $\mathcal O(dL_{-\epsilon,s})$ 的一个小截面，我们希望把 \term{Faltings 复形}{Faltings complex}{complexe de Faltings}延拓到 $A$ 在\term{整数环}{ring of integers}{anneau des entiers}上的模型 $\mathscr A$，然后应用 Siegel 引理。但讨论 $(a\cdot\operatorname{pr}_i+b\cdot\operatorname{pr}_j)^*\mathcal O(L)$ 的截面的算术性质时必须谨慎：态射 $a\cdot\operatorname{pr}_i+b\cdot\operatorname{pr}_j$ 通常不能延拓为算术概形之间的态射。因此，我们使用\term{立方体定理}{theorem of the cube}{théorème du cube}，将这样的截面转化为 $\mathcal O(a^2\operatorname{pr}_i^*L+b^2\operatorname{pr}_j^*L+ab\mathscr P_{ij})$ 的截面；相应\term{同构}{isomorphism}{isomorphisme}须在 $\Spec R$ 上延拓，使其在\term{零截面}{zero section}{section nulle}上也给出\term{同构}{isomorphism}{isomorphisme}。这里与 §9 一样，$\mathscr P_{ij}$ 表示 Poincaré 类
\[
\mathscr P_{ij}:=(\operatorname{pr}_i+\operatorname{pr}_j)^*L-\operatorname{pr}_i^*L-\operatorname{pr}_j^*L.
\]
于是 $(a\cdot\operatorname{pr}_i+b\cdot\operatorname{pr}_j)^*\mathcal O(L)$ 的截面的整性，应通过相应\term{同构}{isomorphism}{isomorphisme}
\begin{equation*}\tag{13.1}
(a\cdot\operatorname{pr}_i+b\cdot\operatorname{pr}_j)^*\mathcal O(L)
\cong\mathcal O(a^2\operatorname{pr}_i^*L+b^2\operatorname{pr}_j^*L+ab\mathscr P_{ij})
\end{equation*}
来理解。这里及下文，以 $\mathscr W_0$ 表示 $A^n$ 的某个模型，使所有 $\mathscr P_{ij}$ 都定义为其上的 Cartier \term{除子类}{divisor class}{classe de diviseurs}。
\begin{statement}{引理 13.2（[F 1]，引理 5.1）}
固定 $A$ 在 $B$ 上的一个模型 $\mathscr A$、$\mathscr A$ 上的一个\term{算术除子}{arithmetic divisor}{diviseur arithmétique}类 $L$，以及在\term{一般纤维}{generic fiber}{fibre générique} $A$ 上生成 $\mathcal O(L)$ 的截面 $\gamma_1,\ldots,\gamma_m\in\Gamma(\mathscr A,\mathcal O(L))$。设 $\mathscr A^2$ 是 $A\times A$ 的某个模型（如上）。则存在常数 $c_{5,v}$，几乎所有 $v$ 处都有 $c_{5,v}=0$，并满足以下性质。
设 $a,b\in\mathbb Z$ 不同时为零。对所有\term{阿基米德赋值}{archimedean place}{place archimédienne} $v$，截面
\[
(a\cdot\operatorname{pr}_1+b\cdot\operatorname{pr}_2)^*\gamma_i
\in\mathcal O(a^2\operatorname{pr}_1^*L+b^2\operatorname{pr}_2^*L+ab\mathscr P_{12})
\]
满足：
\begin{enumerate}[label=(\alph*).]
\item 对所有 $i=1,\ldots,m$ 和所有 $P\in A^2(\bar k_v)$，
\[
\log\|((a\cdot\operatorname{pr}_1+b\cdot\operatorname{pr}_2)^*\gamma_i)(P)\|_v\le c_{5,v}(a^2+b^2).
\]
\item 对所有 $P\in A^2(\bar k_v)$，
\[
\max_{1\le i\le m}\log\|((a\cdot\operatorname{pr}_1+b\cdot\operatorname{pr}_2)^*\gamma_i)(P)\|_v\ge-c_{5,v}(a^2+b^2).
\]
\end{enumerate}
\origpage{196}{191}
再设 $X$ 是 $A^2$ 的一个闭整子概形，$\mathscr X$ 是其在 $\mathscr A^2$ 中的闭包。设 $v$ 是\term{非阿基米德赋值}{non-archimedean place}{place non archimédienne}，$F$ 是 $\mathscr X$ 在 $v$ 处纤维的一个分支，$e$ 是该分支在此纤维中的重数。则对所有这样的 $\mathscr X,v,F$，均有：
\begin{enumerate}[label=(\alph*).,start=3]
\item 对所有 $i=1,\ldots,m$，除子 $((a\cdot\operatorname{pr}_1+b\cdot\operatorname{pr}_2)^*\gamma_i)$ 在 $F$ 处的重数不小于 $-c_{5,v}e(a^2+b^2)$。
\item 当 $i=1,\ldots,m$ 变化时，除子 $((a\cdot\operatorname{pr}_1+b\cdot\operatorname{pr}_2)^*\gamma_i)$ 在 $F$ 处重数的最小值至多为 $c_{5,v}e(a^2+b^2)$。
\end{enumerate}
\end{statement}
\begin{proof}
首先注意，若 $v$ 是好约化赋值，则由纤维上的\term{立方体定理}{theorem of the cube}{théorème du cube}，(c)、(d) 可取 $c_{5,v}=0$。还应注意，这些常数与 $X$ 无关。

我们给出关于 \term{Weil 函数}{Weil functions}{fonctions de Weil}的两个断言，它们分别同时推出 (a)、(c) 和 (b)、(d)，从而证明所需结果。关于 \term{Weil 函数}{Weil function}{fonction de Weil}，参见（[L 5]，第十章）。具体地，对 $i=1,\ldots,m$ 令 $D_i=((a\cdot\operatorname{pr}_1+b\cdot\operatorname{pr}_2)^*\gamma_i)$。则 (a)、(c) 来自不等式
\begin{equation*}\tag{13.3}
\lambda_{D_i}(P)\ge-c_{5,v}(a^2+b^2),
\quad i=1,\ldots,m,\quad P\in A^2(\bar k_v),
\end{equation*}
而 (b)、(d) 来自
\begin{equation*}\tag{13.4}
\min_{1\le i\le m}\lambda_{D_i}(P)\le c_{5,v}(a^2+b^2),
\qquad P\in A^2(\bar k_v).
\end{equation*}
如 (13.1) 中一样，这些 \term{Weil 函数}{Weil function}{fonction de Weil}应由 $L$ 和 $\mathscr P$ 的某些表示上的固定 \term{Weil 函数}{Weil function}{fonction de Weil}，借助\term{立方体定理}{theorem of the cube}{théorème du cube}导出。非阿基米德情形 (c)、(d) 可由上述断言推出：只需选取 $P\in X(\bar k_v)$，使其截面与 $F$ 横截相交。

为证明这两个断言，我们利用下述结果，选取称为 \term{Néron 函数}{Néron functions}{fonctions de Néron}的典范 \term{Weil 函数}{Weil function}{fonction de Weil}。
\begin{statement}{定理 13.5}
设 $A$ 是定义在 $k_v$ 上的\term{阿贝尔簇}{abelian variety}{variété abélienne}。对 $A$ 上支集不含 $0$ 的每个除子 $D$，存在一个 \term{Weil 函数}{Weil function}{fonction de Weil} $\hat\lambda_D$，由以下性质唯一确定：
\begin{enumerate}[label=(\arabic*).]
\item $\hat\lambda_D(0)=0$。
\item 若 $D,D'$ 是 $A$ 上支集不含 $0$ 的除子，则
\[
\hat\lambda_{D+D'}=\hat\lambda_D+\hat\lambda_{D'}.
\]
\item 若 $D=(f)$ 是\term{主除子}{principal divisor}{diviseur principal}，则 $\hat\lambda_D(P)=v(f(P)/f(0))$。
\item $\hat\lambda_{[2]^*D}=\hat\lambda_D\circ[2]$。
\end{enumerate}
这些函数还满足：若 $\phi:B\to A$ 是定义在 $k_v$ 上的\term{阿贝尔簇}{abelian variety}{variété abélienne}同态，则
\begin{equation*}\tag{13.6}
\hat\lambda_{\phi^*D}=\hat\lambda_D\circ\phi.
\end{equation*}
\end{statement}
\begin{proof}
参见（[L 5]，第十一章，定理 1.1）。归一化条件 (1) 消去了该处使用的常数。
\end{proof}
\origpage{197}{192}
对每个固定除子 $D$，$\lambda_D$ 与 $\hat\lambda_D$ 之差的\term{绝对值}{absolute value}{valeur absolue}有界，所以只需对 $\hat\lambda$ 证明 (13.3)、(13.4)。这要求 $D_i$ 不经过 $0$；必要时先按 $A(k_v)$ 中某个固定点作平移，即可保证这一点。

由 (13.6)，无论通过 (13.1) 计算，还是直接通过 $(a\cdot\operatorname{pr}_i+b\cdot\operatorname{pr}_j)^*$ 的定义计算，所得 $\hat\lambda_{D_i}$ 均相同。由于所有除子 $(\gamma_i)$ 都有效，可以选取 $c_{5,v}$，使对所有 $P$ 和所有 $i$ 均有
\[
\hat\lambda_{(\gamma_i)}(P)\ge-c_{5,v},
\]
这就推出 (13.3)。又因这些除子在\term{一般纤维}{generic fiber}{fibre générique}上没有共同的几何点，可以选取 $c_{5,v}$，使对所有 $P$ 均有
\[
\min_{1\le i\le m}\hat\lambda_{(\gamma_i)}(P)\le c_{5,v},
\]
从而推出 (13.4)。
\end{proof}

这一结果的一个较弱形式也在（[V 5]，引理 6.5）中得到证明。但这里需要更强的断言，因为 (c)、(d) 中的子概形 $X$ 可能变化。

将该引理应用于 \term{Faltings 复形}{Faltings complex}{complexe de Faltings} (9.4)，由引理 13.2 的 (c)、(d) 得到以下推论。
\begin{statement}{推论 13.7}
将各映射乘以大小至多为 $\exp(cd\sum s_i^2)$ 的整数后，\term{Faltings 复形}{Faltings complex}{complexe de Faltings}可延拓成一个 $R$ 模序列；所得复形的上同调为\term{挠模}{torsion module}{module de torsion}，并被具有同类上界的整数零化。
\end{statement}
这里再给出两个引理，下一节也会用到。
\begin{statement}{引理 13.8（[F 1]）}
设 $X_1,\ldots,X_n$ 是 $X$ 的子簇，$\mathscr X_1,\ldots,\mathscr X_n$ 是其在 $\mathscr A$ 中的闭包。则对所有充分大的自然数 $d_1,\ldots,d_n$，
\[
V\!\left(\Gamma\!\left(\prod\mathscr X_i,\sum d_i\operatorname{pr}_i^*L\right)\right)
\ge\exp\!\left(-c\,h^0\!\left(\prod X_i,\sum d_i\operatorname{pr}_i^*L\right)\sum d_i\right),
\]
其中常数 $c$ 不依赖于 $X_1,\ldots,X_n$ 及 $d_1,\ldots,d_n$。
\end{statement}
\begin{proof}
利用 $L$ 的一个\term{很丰沛}{very ample}{très ample}倍数，将 $\mathscr A$ 嵌入射影空间。于是直和 $\bigoplus_{d\ge0}\Gamma(\mathscr A,dL)$ 是齐次多项式环上的有限生成模。因此
\[
V(\Gamma(\mathscr A,dL))\ge\exp(-c\,h^0(X,dL)\,d).
\]
由于 $L$ 在\term{一般纤维}{generic fiber}{fibre générique}上\term{丰沛}{ample}{ample}，限制到 $\mathscr X_i$ 不会使这个 $c$ 增大。又因为 $\Gamma(\prod\mathscr X_i,\sum d_i\operatorname{pr}_i^*L)$ 是这些模的张量积的直和，引理随即成立。
\end{proof}
\begin{statement}{引理 13.9}
设 $X_1,\ldots,X_n$、$\mathscr X_1,\ldots,\mathscr X_n$ 和 $d_1,\ldots,d_n$ 如上，并假设 $X_i$ 的次数有界。则对所有充分大的 $d_1,\ldots,d_n$ 及所有非\term{零截面}{zero section}{section nulle}
$\gamma\in\Gamma(\prod\mathscr X_i,\mathcal O(\sum d_i\operatorname{pr}_i^*L))$，
\[
\prod_{v\mid\infty}\|\gamma\|_{\sup,v}\ge\exp\!\left(-c\sum d_i h(X_i)-c'\sum d_i\right),
\]
其中常数 $c,c'$ 不依赖于 $\gamma$、$X_1,\ldots,X_n$ 和 $d_1,\ldots,d_n$。
\end{statement}
\origpage{198}{193}
\begin{proof}
与下文 (18.5) 中一样，存在次数为 $N_i$ 的投影
$\pi_i:\mathscr X_i\to\mathbb P^{m_i}_{\Spec R}$，其中 $m_i=\dim X_i$，使 $\gamma$ 的范数成为 $\mathcal O(e_1,\ldots,e_n)$ 的一个整截面 $\gamma'$\translatornote{24}，且
\begin{equation*}\tag{13.10}
e_i=d_i\prod_{\substack{j=1\\j\ne i}}^nN_j.
\end{equation*}
无穷处的范数由 \term{Fubini--Study 度量}{Fubini--Study metric}{métrique de Fubini--Study}导出，并且对 $k$ 的所有\term{阿基米德赋值}{archimedean place}{place archimédienne} $v$，及使 $\pi_1\times\cdots\times\pi_n$ 在其上有限的所有 $P\in(\mathbb P^{m_1}\times\cdots\times\mathbb P^{m_n})(\bar k_v)$，均有
\begin{equation*}\tag{13.11}
\prod_{Q\in(\pi_1\times\cdots\times\pi_n)^{-1}(P)}\|\gamma(Q)\|
\ge\|\gamma'(P)\|\exp\!\left(c\sum d_i h(X_i)+c''\sum d_i\right).
\end{equation*}
$\mathbb P^m$ 中齐次坐标全是单位根的点组成一个 \term{Zariski 稠密集}{Zariski dense set}{ensemble dense pour la topologie de Zariski}。因此在 (13.11) 中取 $P=(P_1,\ldots,P_n)$，使每个 $P_i$ 都具有这种形式，并以
\[
E\subseteq\mathbb P^{m_1}_{\Spec R}\times_{\Spec R}\cdots\times_{\Spec R}\mathbb P^{m_n}_{\Spec R}
\]
表示对应于 $P$ 的\term{算术曲线}{arithmetic curve}{courbe arithmétique}。于是
\begingroup\small
\begin{align*}
\frac1{[k:\mathbb Q]}\sum_{v\mid\infty}-\log\|\gamma\|_{\sup,v}
&\le\frac1{[k(P):\mathbb Q]}\sum_{w\mid\infty}-\log\|\gamma'(P)\|_w
+c\sum d_i h(X_i)+c''\sum d_i\\
&\le\frac1{[k(P):\mathbb Q]}\sum_{w\mid\infty}\deg\mathcal O(e_1,\ldots,e_n)|_E
+c\sum d_i h(X_i)+c''\sum d_i\\
&\le\sum_{i=1}^ne_i(h(P_i)+c''')+c\sum d_i h(X_i)+c''\sum d_i\\
&\le c\sum d_i h(X_i)+c'\sum d_i.
\end{align*}
\endgroup
最后一步来自 (13.10)，以及所有 $P_i$ 的朴素高度都为零这一事实。这里 $w$ 遍历 $k(P)$ 的\term{阿基米德赋值}{archimedean place}{place archimédienne}。
\end{proof}
还应指出，应用此引理时，可以使用 (18.5) 的完整版本（包括 Poincaré 除子），然后在 \term{Faltings 复形}{Faltings complex}{complexe de Faltings}中以 (9.3) 代替 (9.4)。

\section*{§14. 整体截面的构造}
\addcontentsline{toc}{section}{§14. 整体截面的构造}
前面三节已经备齐构造一个小\term{整体截面}{global section}{section globale}所需的全部工具。
\origpage{199}{194}
\begin{statement}{命题 14.1}
设 $(s)$ 与 $X_1,\ldots,X_n$ 满足条件 (10.5)，$\epsilon$ 如命题 11.5。以 $\mathscr W$ 表示 $X_1\times\cdots\times X_n$ 在 $\mathscr W_0$ 中的闭包。则对所有充分大的 $d\in\mathbb N$（可依赖于 $s$），存在一个整截面
\[
\gamma\in\Gamma(\mathscr W,dL_{-\epsilon,s}),
\]
使得
\begin{equation*}\tag{14.2}
\prod_{v\mid\infty}\|\gamma\|_{\sup,v}\le\exp\!\left(cd\sum s_i^2\right).
\end{equation*}
\end{statement}
\begin{proof}
令 $\beta:\Gamma_1\to\Gamma_2$ 为 \term{Faltings 复形}{Faltings complex}{complexe de Faltings} (9.4) 中最后一个箭头。现在已将整个构造延拓到 $\Spec R$ 上的概形；例如\translatornote{25}，
\[
\Gamma_1=\Gamma\!\left(\mathscr W,d(2n-2+\epsilon)\sum s_i^2\operatorname{pr}_i^*L\right)^a.
\]
其度量取各分量上确界范数中的最大值。又令
\[
\delta_0=\rank(\Ker\beta),\qquad
\delta_1=\rank(\Gamma_1),\qquad
\delta_2=\rank(\im\beta).
\]
对 $\beta$ 应用引理 12.1，即可完成证明。

首先需要在引理 13.8、13.9 中以 $\mathscr W$ 代替 $\prod\mathscr X_i$。这很容易做到，因为对 $\prod\mathscr X_i$ 上任一固定\term{除子类}{divisor class}{classe de diviseurs} $F$，存在从 $\Gamma(\prod\mathscr X_i,dF)$ 到 $\Gamma(\mathscr W,dF)$ 的单射，其余核被一个与 $d$ 无关的整数零化。因此，由引理 13.8，
\begin{equation*}\tag{14.3}
V(\Gamma_1)\ge\exp\!\left(-\delta_1cd\sum s_i^2\right).
\end{equation*}
由引理 13.9、(10.5.4) 及 (10.9)，对所有 $\gamma\notin\Ker\beta$，
\[
\prod_{v\mid\infty}\|\beta(\gamma)\|\ge\exp\!\left(-cd\sum s_i^2\right),
\]
因而
\[
V(\im\beta)\le\exp\!\left(cd\delta_2\sum s_i^2\right).
\]
将此式与 (14.3) 及引理 12.1(b) 合并，得到
\[
V(\Ker\beta)\ge\exp\!\left(-cd\delta_1\sum s_i^2\right).
\]
显然 $\delta_1\le cd^{\sum\dim X_i}\prod s_i^{2\dim X_i}$。因此，命题 11.5 表明 $\delta_1/\delta_0$ 有界；继而引理 12.1(a) 给出一个截面 $\gamma\in\Ker\beta$，使
\begin{equation*}\tag{14.4}
\prod_{v\mid\infty}\|\gamma\|_v\le\exp\!\left(cd\sum s_i^2\right).
\end{equation*}
这里用引理 13.2(a) 得到 (12.2)。再由引理 13.2(b)，当把 $\gamma$ 看作 $\Gamma(\mathscr W,dL_{-\epsilon,s})$ 中的元素时，(14.4) 也给出其上确界范数的同一上界。
\end{proof}
\origpage{200}{195}
\section*{§15. 若干分析结果}
\addcontentsline{toc}{section}{§15. 若干分析结果}
本节收集下一节处理偏导数所需的若干分析结果。由于技术上的困难，必须在 $k$ 的所有赋值处采用引理 6.2 的分析方法。

对\term{阿基米德赋值}{archimedean place}{place archimédienne} $v$，令 $\mathbb C_v=\bar k_v$；取\term{代数闭包}{algebraic closure}{clôture algébrique}后，该域仍然完备。\term{非阿基米德赋值}{non-archimedean place}{place non archimédienne}处则不然：只有再取 $\bar k_v$ 的\term{完备化}{completion}{complété} $\mathbb C_v$，才能得到既完备又代数闭的域。在这样的域上，\term{幂级数}{power series}{série entière}所定义的函数往往具有类似于 $\mathbb C^m$ 上函数的性质，但也存在差异。例如，$\mathbb C_v$ 是\term{完全不连通}{totally disconnected}{totalement discontinu}的，所以往往需要指明，所讨论的函数在整个定义域上都由同一个\term{幂级数}{power series}{série entière}给出。这套思想通常称为\term{刚性分析}{rigid analysis}{analyse rigide}。关于 $p$ 进分析的更多细节，参见 [Ko]；关于\term{刚性分析}{rigid analysis}{analyse rigide}，参见 [T]，也可从另一角度参见 [Be]。

在非阿基米德情形，\term{绝对值}{absolute value}{valeur absolue} $\|\cdot\|_v$ 可延拓为 $\mathbb C_v$ 上的\term{绝对值}{absolute value}{valeur absolue}；讨论 $\mathbb C_v$ 上的分析时，我们常省去下标 $v$，写成 $|\cdot|$。在阿基米德情形，$|\cdot|$ 表示通常的\term{绝对值}{absolute value}{valeur absolue}。对向量 $x=(x_1,\ldots,x_m)\in\mathbb C_v^m$，阿基米德情形令 $|x|^2=|x_1|^2+\cdots+|x_m|^2$，非阿基米德情形则令 $|x|=\max(|x_1|,\ldots,|x_m|)$。以 $D_\rho$ 和 $\bar D_\rho$ 分别表示开圆盘 $\{z\in\mathbb C_v^m:|z|<\rho\}$ 和闭圆盘 $\{z\in\mathbb C_v^m:|z|\le\rho\}$。这些记号沿用自阿基米德情形，稍有误导之处，因为非阿基米德情形下这两个集合都既开又闭。另须注意，非阿基米德情形下，\term{多圆盘}{polydisc}{polydisque}与球并无区别。

\term{非阿基米德赋值}{non-archimedean place}{place non archimédienne}处，Cauchy 不等式仍然成立。
\begin{statement}{引理 15.1}
设 $v$ 是\term{非阿基米德赋值}{non-archimedean place}{place non archimédienne}，多元\term{幂级数}{power series}{série entière}
\[
f(z)=\sum_{(i)\ge0}a_{(i)}z^{(i)}
\]
在半径 $\rho>0$ 的开圆盘 $D_\rho$ 中绝对收敛。则对所有元组 $(\ell)\ge0$，
\[
|a_{(\ell)}|\rho^{\ell_1+\cdots+\ell_m}\le\sup_{z\in D_\rho}|f(z)|.
\]
\end{statement}
\begin{proof}
与阿基米德情形一样，只需证明一元情形。又只需对每个 $\rho'<\rho$ 证明 $|a_\ell|(\rho')^\ell\le\sup_{z\in\bar D_{\rho'}}|f(z)|$，然后取极限即可。经伸缩可设 $\rho'=1$；此时\term{幂级数}{power series}{série entière}在 $z=1$ 收敛，故 $|a_i|\to0$。乘以适当常数后可设 $\max_i|a_i|=1$，于是只需在闭单位圆盘中找到一个 $z$，使 $|f(z)|=1$。由 \term{Newton 多边形}{Newton polygon}{polygone de Newton}理论（[Ko]，第四章，§3）及 Weierstrass 预备定理，圆盘边界上任何一个不靠近相应 Weierstrass 多项式根的 $z$ 都满足这一等式。
\end{proof}
我们还将使用 Newton 方法及其 $p$ 进推论——\term{Hensel 引理}{Hensel's lemma}{lemme de Hensel}。这里遵循 Lang（[L 4]，第二章，§2，命题 2）的方法，但在多元情形下论证。首先设 $K$ 是具有非阿基米德\term{绝对值}{absolute value}{valeur absolue} $|\cdot|$ 的域，并令 $A:=\{x\in K:|x|\le1\}$ 为其\term{赋值环}{valuation ring}{anneau de valuation}。注意，此赋值不必离散。\translatornote{26}
\origpage{201}{196}
对向量 $x\in K^m$，仍按上面定义 $|x|$；对 $m\times m$ 矩阵 $J$，令
\[
|J|=\min_{x\ne0}\frac{|Jx|}{|x|}.
\]
若 $J$ 的元素都属于 $A$，则 $|J|\ge|\det J|$，因此下面给出的 \term{Hensel 引理}{Hensel's lemma}{lemme de Hensel}加强了（[G]，命题 5.20）中的版本。又若 $|J|\ne0$，则 $J$ 非奇异，且
\begin{equation*}\tag{15.2}
|J^{-1}x|\le\frac{|x|}{|J|},\qquad x\in K^m.
\end{equation*}
以下就是多元 \term{Hensel 引理}{Hensel's lemma}{lemme de Hensel}。
\begin{statement}{引理 15.3}
设 $\mathbf f=(f_1,\ldots,f_m)\in A[X_1,\ldots,X_m]^m$ 是一个多项式向量，并以\translatornoteref{26}
\[
J_f:=\left(\frac{\partial f_i}{\partial X_j}\right)_{i,j}
\]
表示其 Jacobian 矩阵。设 $\alpha_0\in A^m$ 满足
\begin{equation*}\tag{15.4}
|\mathbf f(\alpha_0)|<|J_f(\alpha_0)|^2.
\end{equation*}
则 $\mathbf f(X)$ 在 $A^m$ 中存在唯一一个根 $\alpha$ 满足
\begin{equation*}\tag{15.5}
|\alpha-\alpha_0|\le\frac{|\mathbf f(\alpha_0)|}{|J_f(\alpha_0)|}<1,
\end{equation*}
且序列
\begin{equation*}\tag{15.6}
\alpha_{i+1}=\alpha_i-J_f(\alpha_i)^{-1}\cdot\mathbf f(\alpha_i)
\end{equation*}
收敛到 $\alpha$。
\end{statement}
\begin{proof}
先考虑满足 (15.5) 的 $\alpha$，显然 $\alpha\in A^m$。我们还要说明 $|J_f(\alpha)|=|J_f(\alpha_0)|$。事实上，对 $|x|=1$ 的 $x\in K^m$，由 $J_f$ 的 Taylor 展开，
\[
J_f(\alpha)\cdot x=J_f(\alpha_0)\cdot x+\beta\cdot(\alpha-\alpha_0),
\]
其中 $\beta$ 是元素属于 $A$ 的某个矩阵（依赖于 $x$）。由 (15.5)、(15.4)，
\[
|\beta\cdot(\alpha-\alpha_0)|<|J_f(\alpha_0)|.
\]
这在任何情形下都推出 $|J_f(\alpha)|\le|J_f(\alpha_0)|$；取使 $|J_f(\alpha_0)\cdot x|$ 最小的 $x$，又得到反向不等式。

为证明唯一性，设 $\alpha,\alpha'$ 是满足 (15.5) 的两个根。在 $\alpha$ 处对 $\mathbf f$ 作 Taylor 展开，得到
\[
\mathbf f(\alpha')=\mathbf f(\alpha)-J_f(\alpha)\cdot(\alpha'-\alpha)
+{}^t(\alpha'-\alpha)\cdot\beta\cdot(\alpha'-\alpha),
\]
其中 $\beta$ 是由 $m$ 个矩阵组成的元组，各矩阵元素均在 $A$ 中。但 $\mathbf f(\alpha')=\mathbf f(\alpha)=0$，由刚证得的结论，右边第二项的\term{绝对值}{absolute value}{valeur absolue}至少是 $|J_f(\alpha_0)||\alpha-\alpha'|$，而末项的\term{绝对值}{absolute value}{valeur absolue}至多是 $|\alpha-\alpha'|^2$。除非 $\alpha=\alpha'$，否则这与 (15.5)、(15.4) 矛盾。
\origpage{202}{197}
最后令 $c=|\mathbf f(\alpha_0)|/|J_f(\alpha_0)|^2<1$。我们归纳证明：
\begin{enumerate}[label=(\roman*).]
\item $\alpha_i$ 满足 (15.5)；
\item $\displaystyle\frac{|\mathbf f(\alpha_i)|}{|J_f(\alpha_i)|^2}\le c^{2^i}$。
\end{enumerate}
这两个条件显然推出引理：由 (ii)，$\alpha_i$ 收敛到一个根；由 (i)，这个根满足 (15.5)。

$i=0$ 时显然成立。假设对 $i$ 成立。为证明 (i)，注意
\[
|\alpha_{i+1}-\alpha_i|\le c^{2^i}|J_f(\alpha_i)|
\le c|J_f(\alpha_0)|=\frac{|\mathbf f(\alpha_0)|}{|J_f(\alpha_0)|}.
\]
故 (15.5) 成立。证明 (ii) 时再次利用 $\mathbf f$ 的 Taylor 展开：
\begin{align*}
\mathbf f(\alpha_{i+1})
&=\mathbf f(\alpha_i)-J_f(\alpha_i)\cdot J_f(\alpha_i)^{-1}\cdot\mathbf f(\alpha_i)\\
&\quad+{}^t(J_f(\alpha_i)^{-1}\cdot\mathbf f(\alpha_i))\cdot\beta\cdot
(J_f(\alpha_i)^{-1}\cdot\mathbf f(\alpha_i)),
\end{align*}
其中 $\beta$ 仍是由 $m$ 个元素在 $A$ 中的矩阵组成的元组。右边前两项抵消，由 (15.2)，
\begin{equation*}\tag{15.7}
|\mathbf f(\alpha_{i+1})|\le\left(\frac{|\mathbf f(\alpha_i)|}{|J_f(\alpha_i)|}\right)^2,
\end{equation*}
从而得到 (ii)。
\end{proof}
该引理的阿基米德类似形式稍复杂一些。
\begin{statement}{引理 15.8}
设 $\mathbf f=(f_1,\ldots,f_m)$ 是由 $\mathbb C^m$ 到 $\mathbb C$ 的 $C^2$ 函数组成的向量\translatornote{27}，$J_f$ 是如引理 15.3 中的 Jacobian。设 $\alpha_0\in\mathbb C^m$、$B\in\mathbb R_{>0}$ 满足
\begin{equation*}\tag{15.9}
|\mathbf f(\alpha_0)|<\frac{|J_f(\alpha_0)|^2}{2B},
\end{equation*}
以及
\begin{equation*}\tag{15.10}
B>\sup\left|\sum_{i=1}^m\sum_{j=1}^m
\frac{\partial^2\mathbf f(\alpha)}{\partial z_i\partial z_j}v_iw_j\right|,
\end{equation*}
其中 $\alpha$ 遍历以 $\alpha_0$ 为中心、半径为 $\rho$ 的球 $B_\rho(\alpha_0)$，$v,w$ 遍历 $\mathbb C^m$ 中的单位球；这里
\[
\rho=\frac{|\mathbf f(\alpha_0)|}{|J_f(\alpha_0)|}
\sum_{i=0}^{\infty}(1+c)^i\left(\frac{c}{2(1-c)^2}\right)^{2^i-1},
\qquad c=\frac{B|\mathbf f(\alpha_0)|}{|J_f(\alpha_0)|^2}.
\]
则序列 (15.6) 收敛到 $\mathbf f(X)$ 的一个根 $\alpha$，且
\begin{equation*}\tag{15.11}
|\alpha-\alpha_0|\le\frac{|\mathbf f(\alpha_0)|}{|J_f(\alpha_0)|}
\sum_{i=0}^{\infty}(1+c)^i\left(\frac{c}{2(1-c)^2}\right)^{2^i-1}.
\end{equation*}
\end{statement}
\origpage{203}{198}
\begin{proof}
首先注意，若 $\alpha,\alpha'\in B_\rho(\alpha_0)$，则
\begin{equation*}\tag{15.12}
|J_f(\alpha)|-B|\alpha'-\alpha|\le|J_f(\alpha')|
\le|J_f(\alpha)|+B|\alpha'-\alpha|.
\end{equation*}
为证明此式，取 $|x|=1$ 且 $|J_f(\alpha)x|=|J_f(\alpha)|$ 的 $x$。对 $J_f\cdot x$ 应用 Taylor 公式及 (15.10)，得到
\begin{align*}
|J_f(\alpha')|&\le|J_f(\alpha')x|\\
&\le|J_f(\alpha)x|+B|\alpha'-\alpha|\\
&\le|J_f(\alpha)|+B|\alpha'-\alpha|.
\end{align*}
这证明了后一半不等式；前一半由对称性得到。

接下来与非阿基米德情形一样，归纳证明：
\begin{enumerate}[label=(\roman*).]
\item $\displaystyle |\alpha_i-\alpha_0|\le\frac{|\mathbf f(\alpha_0)|}{|J_f(\alpha_0)|}
\sum_{j=1}^{i-1}(1+c)^j\left(\frac{c}{2(1-c)^2}\right)^{2^j-1}$；
\item $|J_f(\alpha_i)|\le(1+c)^i|J_f(\alpha_0)|$；
\item $\displaystyle\frac{B|\mathbf f(\alpha_i)|}{|J_f(\alpha_i)|^2}
\le c\left(\frac{c}{2(1-c)^2}\right)^{2^i-1}$。
\end{enumerate}
由 (15.9)，$c<1/2$，故 $c/[2(1-c)^2]<1$。于是 (ii)、(iii) 推出 $\alpha_i$ 收敛到一个根，(i) 则推出 (15.11)。

为进行归纳，先注意 $\alpha_i$ 的 (ii)、(iii) 推出 $\alpha_{i+1}$ 的 (i)：
\begin{align*}
|\alpha_{i+1}-\alpha_i|
&\le\frac{|\mathbf f(\alpha_i)|}{|J_f(\alpha_i)|}\\
&\le\frac{c|J_f(\alpha_0)|}{B}(1+c)^i
\left(\frac{c}{2(1-c)^2}\right)^{2^i-1}\\
&\le\frac{|\mathbf f(\alpha_0)|}{|J_f(\alpha_0)|}(1+c)^i
\left(\frac{c}{2(1-c)^2}\right)^{2^i-1}.
\end{align*}
再由 (15.12) 的后一半及 $\alpha_i$ 的 (iii)，
\[
|J_f(\alpha_{i+1})|\le|J_f(\alpha_i)|+\frac{B|\mathbf f(\alpha_i)|}{|J_f(\alpha_i)|}
\le|J_f(\alpha_i)|(1+c),
\]
即得 (ii)。为证明 (iii)，先与前面一样注意到
\[
|J_f(\alpha_{i+1})|\ge|J_f(\alpha_i)|(1-c).
\]
与非阿基米德情形一样，代替 (15.7)，我们得到
\[
|\mathbf f(\alpha_{i+1})|\le\frac B2
\left(\frac{|\mathbf f(\alpha_i)|}{|J_f(\alpha_i)|}\right)^2.
\]
\origpage{204}{199}
合并这两式，
\begin{align*}
\frac{B|\mathbf f(\alpha_{i+1})|}{|J_f(\alpha_{i+1})|^2}
&\le\frac{B^2|\mathbf f(\alpha_i)|^2}{2(1-c)^2|J_f(\alpha_i)|^4}\\
&\le\frac{c^2}{2(1-c)^2}
\left(\frac{c}{2(1-c)^2}\right)^{2^{i+1}-2}.
\end{align*}
这就推出 (iii)。
\end{proof}
\noindent\textbf{备注。}阿基米德情形下，唯一性断言不如前面简洁。令 $\rho'=|J_f(\alpha_0)|/(2B)$；若 (15.10) 对 $\alpha\in B_{\rho'}(\alpha_0)$ 成立，则这个球内至多有 $\mathbf f$ 的一个零点。另须注意，$c$ 充分小时 $\rho'\ge\rho$。但本文不需要上述任何一种唯一性断言。

本文将把 \term{Hensel 引理}{Hensel's lemma}{lemme de Hensel}应用于\term{有限态射}{finite morphism}{morphisme fini} $q:X\to\mathbb C_v^m$，其中 $X$ 是 $\mathbb C_v^M$ 的\term{解析子簇}{analytic subvariety}{sous-variété analytique}。设 $P$ 是 $q$ 为 étale 的一点；则 $X$ 中 $P$ 的某个邻域，作为 $\mathbb C_v^M$ 的子集，可由一个\term{正则序列}{regular sequence}{suite régulière}
\[
f_1,\ldots,f_{M-m}\in A[X_1,\ldots,X_M]
\]
给出。对 $i=1,\ldots,M-m$ 令 $g_i=f_i$，再令 $g_{M-m+1},\ldots,g_M$ 为 $q$ 的坐标。\term{Hensel 引理}{Hensel's lemma}{lemme de Hensel}给出 $q$ 的逐点逆函数；\term{逆函数定理}{inverse function theorem}{théorème d'inversion locale}进一步表明，该函数局部可由各变量的\term{幂级数}{power series}{série entière}表示。但我们需要一个\term{多圆盘}{polydisc}{polydisque}半径的下界，使此\term{幂级数}{power series}{série entière}在该\term{多圆盘}{polydisc}{polydisque}上既收敛，又给出 $q$ 的局部逆。这只需在另一个域上应用 \term{Hensel 引理}{Hensel's lemma}{lemme de Hensel}，便容易得到。
\begin{statement}{推论 15.13}
令 $J_g$ 为矩阵 $(\partial g_i/\partial x_j)_{i,j}$，并假设 $J_g$ 在一点 $P\in\mathbb C_v^m$ 非奇异。
\begin{enumerate}[label=(\alph*).]
\item 假设 $P\in A^m$。则 $q$ 的逆在以 $q(P)$ 为中心、半径为 $|J_g(P)|^2$ 的\term{多圆盘}{polydisc}{polydisque}上，由同一个\term{幂级数}{power series}{série entière}给出。
\item 设 $\rho,B$ 如引理 15.8。则 $q$ 的逆在以 $q(P)$ 为中心、半径为 $\rho$ 的球上，由同一个\term{幂级数}{power series}{série entière}给出。
\end{enumerate}
\end{statement}
\begin{proof}
在域
\[
\mathbb C_v((X_1,\ldots,X_M))
\]
上应用 \term{Hensel 引理}{Hensel's lemma}{lemme de Hensel}，并以指标作为赋值（各变量的权可任意选取），便得到该\term{幂级数}{power series}{série entière}。还须说明：它在所指明的\term{多圆盘}{polydisc}{polydisque}中收敛，且代入数值得到的正是先前构造的函数。为此，比较这次应用 \term{Hensel 引理}{Hensel's lemma}{lemme de Hensel}所得的 $\alpha_i$，与逐点应用 \term{Hensel 引理}{Hensel's lemma}{lemme de Hensel}所得的 $\alpha_i$。两组逼近项是一致的：把\term{幂级数}{power series}{série entière} $\alpha_i$ 在某个给定点代入求值，所得数值等于先前方法中的相应逼近项。并且，在任何半径严格较小的\term{多圆盘}{polydisc}{polydisque}上，$|\alpha_{i+1}-\alpha_i|$ 的界都是一致的，故该\term{幂级数}{power series}{série entière}在给定\term{多圆盘}{polydisc}{polydisque}上收敛。
\end{proof}
\origpage{205}{200}
\section*{§16. 导数的进一步讨论}
\addcontentsline{toc}{section}{§16. 导数的进一步讨论}
Faltings 给出的高维推广也使用 $X_1\times\cdots\times X_n$ 上的偏导数，但现在每个 $X_i$ 可能具有更高维数，也可能带有奇点。为处理这些奇点，必须采用上一节的 $p$ 进分析。并且，各簇 $X_i$ 可能变化，因此还需要 \term{Chow 坐标}{Chow coordinates}{coordonnées de Chow}理论。为此，令 $C$ 表示 $\mathbb P^m$ 中给定次数的子簇的 \term{Chow 簇}{Chow variety}{variété de Chow}，$\Gamma$ 表示 $C$ 上的泛簇。这里使用由\term{很丰沛}{very ample}{très ample}除子 $L$ 给出的、已经选定的 $X$（在 $k$ 上）到 $\mathbb P^m$ 的嵌入。在有限多个赋值处，它不能延拓成 $\Spec R$ 上的嵌入；这一问题将在后面处理。

假设这些簇 $C$ 都是射影的。$C$ 中不对应于 $\mathbb P^m$ 的整子簇的点构成一个代数子集；可设各 $X_i$ 均不对应于这样的点。否则，将 $C$ 换为次数严格更小的 \term{Chow 簇}{Chow variety}{variété de Chow}即可。又因 $X_i$ 的次数有界，只需有限多个这样的 $C$。在每个 $C$ 上，对应于 $X_i$ 的点的高度与 §10 中定义的 $h(X_i)$ 有关；详见（[So 1]，定理 3）。
\begin{statement}{命题 16.1}
令 $\mathscr L=\mathcal O(dL_{-\epsilon,s})$，并设 $\gamma$ 是 $\mathscr L$ 在 $X_1\times\cdots\times X_n$ 的模型 $\mathscr W$ 上的\term{整体截面}{global section}{section globale}；此模型由在 $A^n$ 的模型 $\mathscr W_0$ 中取闭包得到。设 $P_0=(P_1,\ldots,P_n)$、$E=E_1\times_B\cdots\times_BE_n$ 及 $d_1,\ldots,d_n$ 如推论 6.3。则存在仅依赖于 $\mathscr W_0,X,\deg X_1,\ldots,\deg X_n$ 的常数 $c_1,c_2,c_3,c_4,c_5,c_6$，以及子集 $Z_i\subseteq X_i$，满足对所有 $i$ 均有 $\deg Z_i\le c_1$、$h(Z_i)\le c_2h(X_i)$，并且以下两种情形之一成立：
\begin{enumerate}[label=(\alph*).]
\item 对某个 $i$，$P_i\in Z_i$；
\item 指标 $t=t(\gamma,P_0,d_1,\ldots,d_n)$ 满足
\[
t\ge\frac{-\deg\mathscr L|_E-\sum_{v\mid\infty}\log\|\gamma\|_{\sup,v}
-c_3\sum d_i-c_4\sum d_i h(X_i)}
{c_5\max_{1\le i\le n}d_i h_L(P_i)+c_6\sum d_i}.
\]
\end{enumerate}
\end{statement}
\begin{proof}
回顾一下，对每个 $i=1,\ldots,n$，已经给定 $X_i$ 到 $\mathbb P^m$ 的嵌入。令 $m_i=\dim X_i$，从 $\mathbb P^m$ 到 $\mathbb P^{m_i}$ 的标准投影中选取一个，使其在 $P_i$ 有定义。对 $\mathbb P^m$ 作\term{爆破}{blowing-up}{éclatement}，使投影在整个\term{爆破}{blowing-up}{éclatement}空间上成为态射。令 $Z_i$ 为此投影的分歧子集。显然，$Z_i$ 的次数和高度满足所需条件。因此，从此设 $P_i\notin Z_i$，专门证明 (b)。又经坐标置换，可设 $P_i$ 不位于 $\mathbb P^{m_i}$ 的无穷远超平面上方。这样便得到投影 $q_i:X_i\to\mathbb A^{m_i}$。

现在将这一构造同时延拓到 \term{Chow 簇}{Chow variety}{variété de Chow}上和 $\Spec R$ 上。令 $C_i$ 为适用于 $X_i$ 的 \term{Chow 簇}{Chow variety}{variété de Chow}，并限制到对应于 $X$ 的子簇的子集。类似地，在 $\Spec R$ 上定义 Chow 概形 $\mathscr C_i$。令 $\Gamma_i\subseteq\mathscr C_i\times_{\Spec R}\mathscr X$ 为\term{泛族}{universal family}{famille universelle}；像前面一样作\term{爆破}{blowing-up}{éclatement}，使 $\Gamma_i$ 能投影到 $\mathbb P^{m_i}$，再延拓模型，使 $\mathscr L$ 延拓为 $\Gamma_i$ 上的\term{线层}{line sheaf}{faisceau inversible}；这些步骤可以不依赖于 $s$ 等数据而完成。再延拓 $C_i$ 和 $\mathscr C_i$，使它们分别在 $\Spec k$ 和 $\Spec R$ 上固有。此外，假设 $\Gamma_i$ 在 $\mathscr C_i$ 上射影，即对某个 $M$ 有 $\Gamma_i\subseteq\mathbb P^M_{\mathscr C_i}$。最后，将 $C_i$ 的坏子集定义为对应子簇不整，或不满射到 $\mathbb P^{m_i}$ 的点集。
\origpage{206}{201}
注意，至此作出的选择都不依赖于赋值 $v$。对每个 $v$，$q_i$ 延拓为从 $\mathscr X_i(\mathbb C_v)$ 的一个子集到 $\mathbb C_v^{m_i}$ 的 $v$ 进解析态射。由对 $P_i$ 的假设，它在某个半径为 $\rho_{i,v}>0$ 的开\term{多圆盘}{polydisc}{polydisque}上局部\term{双全纯}{biholomorphic}{biholomorphe}。下一步就是控制这个半径。

稍作模型变换后，将 $\mathscr X_i$ 视为 $\Gamma_i$ 的子集，并令 $E_i$ 为 $\mathscr X_i$ 中对应于 $P_i$ 的\term{算术曲线}{arithmetic curve}{courbe arithmétique}。这不影响定义在\term{一般纤维}{generic fiber}{fibre générique}上的指标。
\begin{statement}{引理 16.2}
对每个 $i$，$\Gamma_i$ 上存在一个仅依赖于 $\mathscr C_i$、到 $\mathbb P^{m_i}$ 的投影的选择，以及 $\mathbb P^{m_i}$ 上的坐标置换的\term{算术除子}{arithmetic divisor}{diviseur arithmétique} $D_i$，使对 $k$ 的所有赋值 $v$，映射 $q_i$ 均将 $P_i(\mathbb C_v)$ 在 $\mathscr X_i(\mathbb C_v)$ 中的某个邻域 $U_{i,v}$ \term{双全纯}{biholomorphic}{biholomorphe}地映到 $\mathbb C_v^{m_i}$ 中以 $q_i(P_i)$ 为中心、半径为 $\exp(-(D_i\cdot E_i)_v)$ 的\term{多圆盘}{polydisc}{polydisque}上。此外，$E_i$ 不包含于 $D_i$ 的支集中。最后，在\term{阿基米德赋值}{archimedean place}{place archimédienne}处，还可给出集合 $U_{i,v}$ 在 $\mathbb A^M$ 中半径的上界。
\end{statement}
\begin{proof}
先考虑\term{非阿基米德赋值}{non-archimedean place}{place non archimédienne} $v$。由于 $P_i\notin Z_i$，$q_i$ 在 $P_i$ 处 étale，故 $\Gamma_i$ 在 $P_i$ 处相对于 $\mathscr C_i$ 光滑（这里考察的是 $\Spec k$ 上的\term{一般纤维}{generic fiber}{fibre générique}）。回顾 $E_{i,v}$ 表示 $E_i$ 中位于 $v\in\Spec R$ 上方的闭点。由光滑性（[F-L]，第四章，命题 3.11），存在\term{正则序列}{regular sequence}{suite régulière} $f_1,\ldots,f_{M-m_i}$，给出 $\Gamma_i$ 中 $E_{i,v}$ 的某个邻域作为某个仿射空间 $\mathbb A^M_{\mathscr C_i}$ 的子集的理想。这些函数以 $\mathbb P^M$ 上的局部坐标 $x_1,\ldots,x_M$ 及 $\mathscr C_i$ 上的局部坐标表示。令 $y_1,\ldots,y_{m_i}$ 为 $q_i$ 的坐标，则 $y_1-y_1(P_i),\ldots,y_{m_i}-y_{m_i}(P_i)$ 也构成 $P_i$ 在 $X_i$ 上的\term{正则序列}{regular sequence}{suite régulière}。现在专化到 $\mathscr X_i(\mathbb C_v)$；这些函数便在以 $P_i(\mathbb C_v)$ 为中心的开单位圆盘中描述 $\mathscr X_i(\mathbb C_v)$。由于
$f_1,\ldots,f_{M-m_i},y_1-y_1(P_i),\ldots,y_{m_i}-y_{m_i}(P_i)$ 是\term{正则序列}{regular sequence}{suite régulière}，其 Jacobian 行列式非零。取半径为此行列式\term{绝对值}{absolute value}{valeur absolue}的平方。由 \term{Hensel 引理}{Hensel's lemma}{lemme de Hensel}（推论 15.13），所需的\term{双全纯}{biholomorphic}{biholomorphe}映射存在。

当点 $E_{i,v}$ 遍历 $\Gamma_i$ 的某个 Zariski 开子集时，由 Zariski 拓扑的拟紧性，只需有限组这样的 $f_1,\ldots,f_{m_i}$。所以 $\Gamma_i$ 上存在一个除子 $D_i$（此时只在代数部分上，尚非\term{算术除子}{arithmetic divisor}{diviseur arithmétique}），支配这些行列式的平方；而且可选取它，使其\term{一般纤维}{generic fiber}{fibre générique}上相对于 $C_i$ 的垂直分支只对应于前面讨论的坏集。

在\term{阿基米德赋值}{archimedean place}{place archimédienne}处，同样可为 $D_i$ 构造一个 \term{Green 函数}{Green's function}{fonction de Green} $g_{D_i}$，使所需\term{双全纯}{biholomorphic}{biholomorphe}映射在半径为 $\exp(-g_{D_i})$ 的\term{多圆盘}{polydisc}{polydisque}上存在。这里利用 $\Gamma_i(\mathbb C_v)$ 的紧性。
\end{proof}
还需要第二个估计，但它容易得多。先在 $\Omega_{\Gamma_i/\mathscr C_i}$ 上选取一个 Hermite 度量。
\begin{statement}{引理 16.3}
对每个 $i$，$\Gamma_i$ 上存在一个仅依赖于 $\mathscr C_i$、到 $\mathbb P^{m_i}$ 的投影的选择及 $\mathbb P^{m_i}$ 上坐标置换的\term{算术除子}{arithmetic divisor}{diviseur arithmétique} $F_i$，使对 $k$ 的所有赋值 $v$，$P_i$ 在 $X_i$ 中的\term{理想层}{ideal sheaf}{faisceau d'idéaux}的元素 $q_i^*y_j-y_j(q_i(P_i))$ 的度量满足
\[
-\log\|q_i^*y_j-y_j(q_i(P_i))\|_v\ge-(F_i\cdot E_i)_v.
\]
\origpage{207}{202}
（\term{非阿基米德赋值}{non-archimedean place}{place non archimédienne}处，上述度量由概形结构给出。又通过（[H 2]，II 8.12），将 $P_i$ 的\term{理想层}{ideal sheaf}{faisceau d'idéaux}单射到 $\Omega_{\Gamma_i/\mathscr C_i}$。）
\end{statement}
\begin{proof}
取一个支配由 $q_i$ 的相对微分给出的挠层的除子 $F_i$ 即可。加上 $q_i^*H_\infty$（其中 $H_\infty$ 是 $\mathbb P^{m_i}$ 的无穷远超平面），就能校正无穷远处的增长。
\end{proof}
对每个赋值 $v$，令 $U_v$ 是引理 16.2 中各邻域 $U_{i,v}$ 的乘积。又令 $\gamma_{0,v}$ 是 $\mathscr L$ 在 $U_v$ 上的局部生成元。应由 $\operatorname{pr}_i^*\mathcal O(L)$ 和 Poincaré 除子 $\mathscr P_{ij}$ 的生成元一致地选取它，从而保证
\begin{equation*}\tag{16.4}
\log\inf_{P\in U_v}\frac{\|\gamma_{0,v}(P)\|}{\|\gamma_{0,v}(P_0)\|_v}
\ge-c_{3,v}\sum_{i=1}^n d_i.
\end{equation*}
在\term{阿基米德赋值}{archimedean place}{place archimédienne}处，由于 $\mathscr L$ 定义在 $X$ 上，且如引理 16.2 的证明后所述，可以限制 $U_v$ 的半径，故能做到这一点。在 $\mathscr A$ 具有好约化的赋值处，此下确界恰好为 $1$。其他赋值处应按阿基米德情形处理，这是因为定义 Poincaré 除子时需要变换模型。此外，对非阿基米德 $v$，令 $\|\gamma\|_{\sup,v}=1$。

现在取元组 $(\ell)=(\ell_{11},\ldots,\ell_{1m_1},\ldots,\ell_{nm_n})$，使
\begin{equation*}\tag{16.5}
\sum_{i=1}^n\sum_{j=1}^{m_i}\frac{\ell_{ij}}{d_i}=t,
\end{equation*}
并使（以 $y_{ij}$，$j=1,\ldots,m_i$ 表示 $\mathbb C_v^{m_i}$ 的坐标）
\[
D_{(\ell)}\gamma:=\frac1{\ell_{11}!\cdots\ell_{nm_n}!}
\left(\frac\partial{\partial y_{11}}\right)^{\ell_{11}}\cdots
\left(\frac\partial{\partial y_{nm_n}}\right)^{\ell_{nm_n}}\gamma\ne0,
\]
其中偏导数通过映射 $q_i$ 来定义。与引理 6.2 的证明一样，由 (16.4)，偏导数的范数满足
\begin{align*}\tag{16.6}
-\log\|D_{(\ell)}\gamma(P_0)\|_v
&\ge-\log\|\gamma\|_{\sup,v}
+\sum_{i=1}^n\sum_{j=1}^{m_i}\ell_{ij}(\log\rho_{i,v}-(F_i\cdot E_i)_v)\\
&\quad+\log\inf_{P\in U_v}\frac{\|\gamma_{0,v}(P)\|}{\|\gamma_{0,v}(P_0)\|_v}\\
&\ge-\log\|\gamma\|_{\sup,v}
+\sum_{i=1}^n\sum_{j=1}^{m_i}\ell_{ij}(\log\rho_{i,v}-(F_i\cdot E_i)_v)
-c_{3,v}\sum_{i=1}^n d_i.
\end{align*}
这里将 $D_{(\ell)}\gamma(P_0)$ 看作向量层
\[
\left.(\mathscr L\otimes S^{\ell_{11}+\cdots+\ell_{1m_1}}\Omega_{\mathscr X_1/\Spec R}
\otimes\cdots\otimes S^{\ell_{n1}+\cdots+\ell_{nm_n}}\Omega_{\mathscr X_n/\Spec R})\right|_E
\]
的一个截面。
\origpage{208}{203}
其中 $E$ 仍为 $\Gamma_1\times_{\Spec R}\cdots\times_{\Spec R}\Gamma_n$ 上对应于 $P_0=(P_1,\ldots,P_n)$ 的\term{算术曲线}{arithmetic curve}{courbe arithmétique}。亦见（[Laf]，§4）。所以，对 $\Gamma_i$ 上某些充分大的\term{线层}{line sheaf}{faisceau inversible} $\mathscr M_i$，有
\[
\sum_v-\log\|D_{(\ell)}\gamma(P_0)\|_v
\le\deg\mathscr L|_E+\sum_{i=1}^n\sum_{j=1}^{m_i}\ell_{ij}\deg\mathscr M_i|_{E_i}.
\]
与 (16.6) 合并，得到
\begin{align*}\tag{16.7}
&\sum_{i=1}^n\sum_{j=1}^{m_i}\ell_{ij}
\left(\deg\mathscr M_i|_{E_i}-\sum_v\log\rho_{i,v}+(F_i\cdot E_i)\right)\\
&\quad\ge-\deg\mathscr L|_E-\sum_{v\mid\infty}\log\|\gamma\|_{\sup,v}-c_3\sum_{i=1}^n d_i.
\end{align*}
左边括号内是 $\Gamma_i$ 上的一个交数。由 $\Gamma_i$ 作为乘积的子集的结构，可在 $\mathbb P^M$ 上取除子 $G_i$，在 $\mathscr C_i$ 上取除子 $H_i$，使相对于 $\Gamma_i$ 在 $\Spec k$ 上的\term{一般纤维}{generic fiber}{fibre générique}上的\term{很丰沛}{very ample}{très ample}\term{线层}{line sheaf}{faisceau inversible}锥，有
\[
\mathscr M_i\otimes\mathcal O(F_i+D_i)\le\mathcal O(G_i+H_i).
\]
因此
\[
\deg\mathscr M_i|_{E_i}+(F_i\cdot E_i)+(D_i\cdot E_i)
\le(G_i\cdot E_i)+(H_i\cdot E_i)+O(1).
\]
又有
\[
(G_i\cdot E_i)\le c_5h_L(P_i)+O(1),\qquad
(H_i\cdot E_i)\le c_4h(X_i)+O(1).
\]
于是由引理 16.2，(16.7) 变为
\begin{align*}
\sum_{i=1}^n\sum_{j=1}^{m_i}\ell_{ij}
(c_5h_L(P_i)+c_4h(X_i)+c_6)
\ge-\deg\mathscr L|_E-\sum_{v\mid\infty}\log\|\gamma\|_{\sup,v}-c_3\sum_{i=1}^n d_i.
\end{align*}
对所有 $i$ 都有 $\sum_{j=1}^{m_i}\ell_{ij}\ll d_i$，故适当修改 $c_4$ 并利用 (16.5)，得到
\begin{align*}
&c_4\sum_{i=1}^n d_i h(X_i)+c_5t\max_{1\le i\le n}d_i h_L(P_i)+c_6t\sum_{i=1}^n d_i\\
&\quad\ge-\deg\mathscr L|_E-\sum_{v\mid\infty}\log\|\gamma\|_{\sup,v}-c_3\sum_{i=1}^n d_i.
\end{align*}
解出 $t$，即得命题。
\end{proof}

\section*{§17. 指标的下界}
\addcontentsline{toc}{section}{§17. 指标的下界}
与 §7 一样，由 Mordell--Weil 定理，$A(k)$ 是有限生成阿贝尔群，故 $A(k)\otimes_{\mathbb Z}\mathbb R$ 是有限维向量空间。以 $L$ 代替 theta 除子，则
\[
(P_1,P_2)_L:=\hat h_L(P_1+P_2)-\hat h_L(P_1)-\hat h_L(P_2)
\]
在 $A(k)\otimes_{\mathbb Z}\mathbb R$ 上定义一个非退化内积。又令 $|P|^2=(P,P)_L$。
\origpage{209}{204}
现假设已选取 $P_1,\ldots,P_n$，使对某个给定的 $\epsilon_1>0$ 及所有 $i<j$ 均有
\begin{align*}\tag{17.1}
(P_i,P_j)_L&\ge(1-\epsilon_1)\sqrt{|P_i|^2|P_j|^2}\\
&\ge2(1-\epsilon_1)\sqrt{\hat h_L(P_i)\hat h_L(P_j)}.
\end{align*}
由\term{典范高度}{canonical height}{hauteur canonique}的基本性质，
\begin{align*}
\deg((s_i^2\cdot\operatorname{pr}_i-s_j^2\cdot\operatorname{pr}_j)^*L)|_E
&=s_i^2\hat h_L(P_i)+s_j^2\hat h_L(P_j)-s_is_j(P_i,P_j)_L+O(s_i^2+s_j^2)\\
&\le\left(s_i\sqrt{\hat h_L(P_i)}-s_j\sqrt{\hat h_L(P_j)}\right)^2\\
&\quad+2\epsilon_1s_is_j\sqrt{\hat h_L(P_i)\hat h_L(P_j)}+O(s_i^2+s_j^2).
\end{align*}
取接近 $1/\sqrt{\hat h_L(P_i)}$ 的有理数 $s_i$，上式中的平方项便趋近于零，因而得到
\begin{equation*}\tag{17.2}
\deg L_{-\epsilon,s}|_E\le n(n-1)\epsilon_1-n\epsilon+O\!\left(\sum s_i^2\right).
\end{equation*}
现在应用命题 16.1。由 (14.2) 和 (10.5.4)，(16.1) 中分式的分子第二项、第四项均以 $c\sum d_i$ 为界；因此由 (17.2)，指标 $t=t(\gamma,(P_1,\ldots,P_n),d_1,\ldots,d_n)$ 满足
\[
t\ge\frac{n(\epsilon-(n-1)\epsilon_1)-c\sum s_i^2}{c_5+c_6\sum s_i^2}.
\]
现假设各高度 $\hat h_L(P_i)$ 充分大，$\epsilon_1$ 充分小，则 $s_i^2$ 很小，上式变为
\begin{equation*}\tag{17.3}
t\ge\epsilon_2,
\end{equation*}
其中 $\epsilon_2>0$ 仅依赖于惯常的数据 $X,\mathscr A,\ldots,(\dim X_1,\ldots,\dim X_n)$。

\section*{§18. 乘积定理}
\addcontentsline{toc}{section}{§18. 乘积定理}
证明的最后一步是像（[F 1]，§6）或 [F 2] 那样应用\term{乘积定理}{product theorem}{théorème du produit}。
\begin{statement}{定理 18.1（[F 1]，§3）}
设 $\Pi=\mathbb P^{m_1}\times\cdots\times\mathbb P^{m_n}$ 是特征零域上的射影空间乘积，并给定 $\epsilon_3>0$。\translatornote{28}则存在数 $r',c_1,c_2,c_3$，满足下述性质。设 $\gamma'$ 是 $\Pi$ 上层 $\mathcal O(e_1,\ldots,e_n)$ 的一个非零\term{整体截面}{global section}{section globale}，在某点 $(x_1,\ldots,x_n)$ 关于 $(e_1,\ldots,e_n)$ 的指标不小于 $\epsilon_3$。若对 $i=1,\ldots,n-1$ 均有 $e_i/e_{i+1}\ge r'$，则存在子簇 $Y_i\subseteq\mathbb P^{m_i}$，它们不全等于相应的 $\mathbb P^{m_i}$，并满足：
\begin{enumerate}[label=(\roman*).]
\item 每个 $Y_i$ 都包含 $x_i$；
\item 各 $Y_i$ 的次数均以 $c_1$ 为界；
\origpage{210}{205}
\item 高度 $h(Y_i)$ 满足不等式
\begin{equation*}\tag{18.2}
\sum e_i h(Y_i)\le c_2\sum_{v\mid\infty}\log\|\gamma'\|_{\sup,v}+c_3\sum_i e_i.
\end{equation*}
\end{enumerate}
\end{statement}
\begin{proof}
参见（[F 1]，定理 3.1 和定理 3.3）。
\end{proof}
注意，此定理推广了 Roth 引理 (4.5)，不过没有给出显式常数。事实上，令 $m_1=\cdots=m_n=1$，则所得 $Y_i$ 中至少有一个必须等于 $x_i$，于是 (18.2) 与 (4.6) 矛盾。

此定理推出一个更适合当前问题的结果。
\begin{statement}{推论 18.3}
设 $X$ 是定义在\term{数域}{number field}{corps de nombres} $k$ 上的射影簇。固定 $X$ 的一个\term{射影嵌入}{projective embedding}{plongement projectif}，并令 $L$ 为 $X$ 上的\term{丰沛}{ample}{ample}除子。设 $X_1,\ldots,X_n$ 是 $X$ 的\term{几何不可约}{geometrically irreducible}{géométriquement irréductible}子簇，均定义在 $k$ 上，其次数（相对于选定的 $X$ 的\term{射影嵌入}{projective embedding}{plongement projectif}）有界。给定 $\epsilon>0$ 和 $\epsilon_2>0$。则存在数 $r,c_1,c_2,c_3,c_4$，满足下述性质。

设 $\gamma$ 是 $\mathcal O(dL_{-\epsilon,s})$ 的一个非零\term{整体截面}{global section}{section globale}\translatornoteref{28}，在点 $(P_1,\ldots,P_n)$ 关于 $(d_1,\ldots,d_n)$（$d_i=ds_i^2$）的指标不小于 $\epsilon_2$，且每个 $P_i\in X_i(k)$。若对所有 $i=1,\ldots,n-1$ 均有 $d_i/d_{i+1}\ge r$，则存在子簇 $X_i'\subseteq X_i$，它们不全等于相应的 $X_i$，且：
\begin{enumerate}[label=(\roman*).]
\item 每个 $X_i'$ 都包含 $P_i$；
\item 每个 $X_i'$ 都\term{几何不可约}{geometrically irreducible}{géométriquement irréductible}，且定义在 $k$ 上；
\item 各 $X_i'$ 的次数均以 $c_1$ 为界；
\item 高度 $h(X_i')$ 满足
\begin{equation*}\tag{18.4}
\sum d_i h(X_i')\le c_2\sum_{v\mid\infty}\log\|\gamma\|_{\sup,v}
+c_3\sum_i d_i h(X_i)+c_4\sum_i d_i.
\end{equation*}
\end{enumerate}
\end{statement}
\begin{proof}
令 $m_i=\dim X_i$。设 $X$ 嵌入于射影空间 $\mathbb P^m$；对每个 $i$，固定从 $\mathbb P^m$ 到 $\mathbb P^{m_i}$ 的一个标准投影，使其在 $X_i$ 上的限制为\term{一般有限有理映射}{generically finite rational map}{application rationnelle génériquement finie}，以 $N_i$ 表示其次数。事实上，$N_i=\deg X_i$，所以 $N_i$ 有界。我们希望从 $\prod X_i$ 向 $\Pi$ 取范数。为此，对每个 $i$，令 $K_i^*$ 为 $K(X_i)$ 的一个有限扩张，并且在 $K(\mathbb P^{m_i})$ 上正规。取 $K_i^*$ 的一个模型 $X_i^*$，使对应于所有 $K(\mathbb P^{m_i})$ 上单射 $K(X_i)\hookrightarrow K_i^*$ 的\term{有理映射}{rational map}{application rationnelle} $X_i^*\to X_i$ 都成为态射。通过这些态射拉回 $\mathcal O(L)$，再取其乘积，得到 $X_i^*$ 上的一个\term{线层}{line sheaf}{faisceau inversible}；它\term{同构}{isomorphism}{isomorphisme}于从 $\mathbb P^{m_i}$ 拉回的 $\mathcal O(1)$ 的某个倍数。像 §16 一样，将其延拓到 $\mathscr C_i/\Spec R$ 上的 \term{Chow 族}{Chow family}{famille de Chow} $\Gamma_i$。此\term{同构}{isomorphism}{isomorphisme}在相差 $\mathscr C_i$ 上的一个除子、$\Spec R$ 上的纤维分支，以及相应的阿基米德处度量变换的意义下成立。这里用 \term{Fubini--Study 度量}{Fubini--Study metric}{métrique de Fubini--Study}赋予 $\mathcal O(1)$ 度量。因此，此\term{同构}{isomorphism}{isomorphisme}的分母可由 $\exp(ch(X_i)+c')$ 控制，相应度量之比也可用同类的界控制。

对 Poincaré 除子 $\mathscr P_{ij}$ 而言，同样的界还不够，因为它们的系数大小为 $s_is_j$。但进一步精细化论证，就能得到所需的界。令 $C_i$ 表示 $\mathscr C_i$ 的\term{一般纤维}{generic fiber}{fibre générique}。则 $\mathscr P_{ij}$ 的范数是
\[
\mathbb P^{m_i}_{\mathscr C_i}\times_{\Spec R}\mathbb P^{m_j}_{\mathscr C_j}
\]
上的一个\term{除子类}{divisor class}{classe de diviseurs}，并且在 $C_i\times C_j$ 的每一点上方都平凡。
\origpage{211}{206}
但对每个点 $Q\in\mathbb P^{m_i}_{C_i}$，此范数在 $\{Q\}\times\mathbb P^{m_j}_{C_j}$ 上的限制代数等价于零，因为它来自 $\mathscr P_{ij}$ 在形如 $\{Q'\}\times X$ 的集合上的限制，其中 $Q'\in\Gamma_i$ 位于 $Q\in X$ 上方。交换两个因子后同样成立。因此，其差是 $C_i\times C_j$ 上 Poincaré 型的除子。从而，范数之比和分母大小的界具有 $\exp(c\sqrt{h(X_i)h(X_j)}+c')$ 的形式。这恰好合适，因为 $\mathscr P_{ij}$ 在 $L_{-\epsilon,s}$ 中的重数为 $s_is_j$，而
\[
s_is_j\sqrt{h(X_i)h(X_j)}\le\frac12\bigl(s_i^2h(X_i)+s_j^2h(X_j)\bigr),
\]
这正是所需的那种界！

于是消去分母后，$\gamma$ 的范数 $\gamma'$ 的度量满足\translatornote{29}
\begin{equation*}\tag{18.5}
\prod_{v\mid\infty}\|\gamma'\|_{\sup,v}
\le\prod_{v\mid\infty}\|\gamma\|_{\sup,v}
\cdot\exp\!\left(c\sum d_i h(X_i)+c'\sum d_i\right).
\end{equation*}
现在计算 $e_i$。首先，映射 $\prod X_i\to\Pi$ 的次数为 $N:=\prod N_i$，故 $e_i=d_iN/N_i$。\translatornoteref{24}由于 $N_i$ 有界，$\gamma'$ 在 $(P_1,\ldots,P_n)$ 下方的点 $(x_1,\ldots,x_n)$ 处的指标有下界，即使将权从 $(d_1,\ldots,d_n)$ 改为 $(e_1,\ldots,e_n)$ 后仍是如此。$N_i$ 有界还意味着存在某个 $r$，使 $d_i/d_{i+1}\ge r$ 推出 $e_i/e_{i+1}\ge r'$。

对 $\gamma'$ 应用定理 18.1，得到子簇 $Y_i\subseteq\mathbb P^{m_i}$。由 (18.2)、(18.5)，
\begin{align*}
\sum d_i h(Y_i)&\le c\sum e_i h(Y_i)\\
&\le c\sum_{v\mid\infty}\log\|\gamma\|_{\sup,v}
+c'\sum d_i h(X_i)+c''\sum d_i+c'''\sum e_i\\
&\le c\sum_{v\mid\infty}\log\|\gamma\|_{\sup,v}
+c'\sum d_i h(X_i)+c''\sum d_i.
\end{align*}
再将这些 $Y_i$ 拉回为 $X_i$ 的子集 $X_i'$。利用 §10 的定义，$h(X_i')$ 可由 $h(Y_i)$ 控制，故条件 (iv) 成立。条件 (i) 由构造成立，(iii) 也容易检查。由于 $Y_i$ 不全等于 $\mathbb P^{m_i}$，$X_i'$ 也不全等于 $X_i$。

只需再保证 (ii) 成立。将 $X_i'$ 与其在 $k$ 上的有限多个（至多 $\dim X_i$ 个）共轭取交，直至包含 $P_i$ 的\term{几何不可约}{geometrically irreducible}{géométriquement irréductible}分支定义在 $k$ 上。以此不可约分支替换 $X_i'$，便得到 (ii)。由于取交次数有界，适当调整常数后，(iii)、(iv) 仍然成立。
\end{proof}
此推论可直接用于定理 0.3 的情形。若 (16.1a) 成立，则很直接地得到归纳步骤；否则 (16.1b) 成立，由 (17.3) 及推论 18.3 得到 $X_i$ 的子簇 $X_i'$，满足 (10.5.1)--(10.5.3)，常数可能需要改变。再由 (14.2)、$X_i$ 的 (10.5.4) 及 (18.4)，$X_i'$ 也满足 (10.5.4)。此外，至少有一个 $X_i'$ 的维数严格小于 $\dim X_i$。至此，定理 0.3 证明的主要部分完成。
\origpage{212}{207}
回到 §10 的总体方案，现在选取
\[
P_1,\ldots,P_n\in X(k)\setminus Z(X)(k),
\]
使其满足：
\begin{enumerate}[label=(\alph*).]
\item 高度 $h_L(P_1)$ 充分大，既与 (10.7) 矛盾，又保证 (17.3) 成立。
\item 对 $i=1,\ldots,n-1$，$h(P_{i+1})/h(P_i)>r''$，其中 $r''$ 是\term{乘积定理}{product theorem}{théorème du produit}的所有应用中出现的 $r$ 的最大值；并设 $r''\ge1$。
\item 对所有可能的 $(\dim X_1,\ldots,\dim X_n)$，条件 (17.1) 都成立。
\end{enumerate}
特别地，$r$、$\epsilon_1$ 以及各个常数 $c$ 都依赖于元组 $(\dim X_1,\ldots,\dim X_n)$，但这样的元组只有有限多个。于是可按 §10 的方案进行归纳，导出矛盾。

至此，定理 0.3 得证。

最后简要说明这一证明与定理 0.2 的证明有何不同。在定理 0.2 的情形，可以证明 $L_{-\epsilon,s}$ \term{丰沛}{ample}{ample}。将 $\epsilon$ 略微减小，就能对 $\mathcal O(dL_{-\epsilon,s})$ 中在 $(P_1,\ldots,P_n)$ 处指标不小于某个适当 $\sigma>0$ 的截面空间的维数给出上界。此上界与 $h^0(X^n,dL_{-\epsilon,s})$ 之间有正的间隔，故更精确的 Siegel 引理形式 (2.4) 允许我们构造一个\term{整体截面}{global section}{section globale} $\gamma$，使它在 $(P_1,\ldots,P_n)$ 处的指标不超过 $\sigma$。这样，第 5 步便融入第 2 步；随后在第 4 步得到矛盾，无需对子簇 $X_i$ 进行归纳。

\section*{参考文献}
\addcontentsline{toc}{section}{第三篇参考文献}
\begin{description}[style=nextline,leftmargin=4em,font=\normalfont\bfseries]
\item[{[Ab]}] S. S. Abhyankar, \emph{Resolution of singularities of arithmetical surfaces}, Arithmetical algebraic geometry (O. F. G. Schilling, ed.), Harper \& Row, New York, 1965, pp. 111--152.
\item[{[Ar]}] M. Artin, \emph{Lipman's proof of resolution of singularities for surfaces}, Arithmetic geometry (G. Cornell and J. H. Silverman, eds.), Springer-Verlag, New York, 1986, pp. 267--287.
\item[{[Be]}] V. Berkovich, \emph{Spectral theory and analytic geometry over non-Archimedean fields}, AMS Surveys and Monographs 33, Amer. Math. Soc., Providence, R.I., 1990.
\item[{[B-G-S]}] J.-B. Bost, H. Gillet, and C. Soulé, \emph{Un analogue arithmétique du théorème de Bezout}, C. R. Acad. Sci. Paris, Sér. I 312 (1991), 845--848.
\item[{[Bo]}] E. Bombieri, \emph{The Mordell conjecture revisited}, Ann. Sc. Norm. Super. Pisa, Cl. Sci., IV 17 (1990), 615--640.
\item[{[B-V]}] E. Bombieri and J. Vaaler, \emph{On Siegel's lemma}, Invent. Math. 73 (1983), 11--32; addendum, Invent. Math. 75 (1984), 377.
\item[{[D]}] F. J. Dyson, \emph{The approximation to algebraic numbers by rationals}, Acta Math. 79 (1947), 225--240.
\item[{[F 1]}] G. Faltings, \emph{Diophantine approximation on abelian varieties}, Ann. Math. 133 (1991), 549--576.
\item[{[F 2]}] G. Faltings, \emph{The general case of S. Lang's conjecture}（待刊）.
\item[{[F-L]}] W. Fulton and S. Lang, \emph{Riemann-Roch algebra}, Grundlehren der mathematischen Wissenschaften 277, Springer-Verlag, New York, 1985.
\item[{[G]}] M. J. Greenberg, \emph{Lectures on forms in many variables}, Mathematics lecture note series, W. A. Benjamin, Inc., New York, 1969.
\item[{[G-S 1]}] H. Gillet and C. Soulé, \emph{Arithmetic intersection theory}, Publ. Math. IHES 72 (1990), 93--174.
\item[{[G-S 2]}] H. Gillet and C. Soulé, \emph{Characteristic classes for algebraic vector bundles with hermitian metric. I}, Ann. Math. 131 (1990), 163--203; II, Ann. Math. 131 (1990), 205--238.
\item[{[G-S 3]}] H. Gillet and C. Soulé, \emph{Analytic torsion and the arithmetic Todd genus}, Topology 30 (1991), 21--54.
\origpage{213}{208}
\item[{[G-S 4]}] H. Gillet and C. Soulé, \emph{Un théorème de Riemann-Roch-Grothendieck arithmétique}, C. R. Acad. Sci. Paris, Sér. I 309 (1989), 929--932.
\item[{[H 1]}] R. Hartshorne, \emph{Ample subvarieties of algebraic varieties}, Lecture Notes in Mathematics 156, Springer-Verlag, New York, 1970.
\item[{[H 2]}] R. Hartshorne, \emph{Algebraic geometry}, Graduate Texts in Mathematics 52, Springer-Verlag, New York, 1977.
\item[{[Ka]}] Y. Kawamata, \emph{On Bloch's conjecture}, Invent. Math. 57 (1980), 97--100.
\item[{[Ko]}] N. Koblitz, \emph{$p$-adic numbers, $p$-adic analysis, and zeta-functions}, Graduate Texts in Mathematics 58, Springer-Verlag, New York, 1977.
\item[{[Laf]}] L. Lafforgue, \emph{Une version en géométrie diophantienne du ``lemma de l'indice''}（待刊）.
\item[{[L 1]}] S. Lang, \emph{Some theorems and conjectures in diophantine equations}, Bull. AMS 66 (1960), 240--249.
\item[{[L 2]}] S. Lang, \emph{Integral points on curves}, Publ. Math. IHES 6 (1960), 27--43.
\item[{[L 3]}] S. Lang, \emph{Introduction to diophantine approximations}, Addison-Wesley, Reading, Mass., 1966.
\item[{[L 4]}] S. Lang, \emph{Algebraic number theory}, Addison-Wesley, Reading, Mass., 1970; reprinted, Springer-Verlag, Berlin-Heidelberg-New York, 1986.
\item[{[L 5]}] S. Lang, \emph{Fundamentals of diophantine geometry}, Springer-Verlag, New York, 1983.
\item[{[L 6]}] S. Lang, \emph{Hyperbolic and diophantine analysis}, Bull. AMS 14 (1986), 159--205.
\item[{[L 7]}] S. Lang, \emph{Introduction to Arakelov theory}, Springer-Verlag, New York, 1988.
\item[{[LeV]}] W. J. LeVeque, \emph{Topics in Number Theory}, Vol. II, Addison-Wesley, Reading, Mass., 1956.
\item[{[M]}] D. Mumford, \emph{A remark on Mordell's conjecture}, Amer. J. Math. 87 (1965), 1007--1016.
\item[{[N]}] J. Noguchi, \emph{A higher dimensional analogue of Mordell's conjecture over function fields}, Math. Ann. 256 (1981), 207--212.
\item[{[Ra]}] M. Raynaud, \emph{Courbes sur une variété abélienne et points de torsion}, Invent. Math. 71 (1983), 207--223.
\item[{[Ro]}] K. F. Roth, \emph{Rational approximations to algebraic numbers}, Mathematika 2 (1955), 1--20; corrigendum, Mathematika 2 (1955), 168.
\item[{[Sch]}] W. M. Schmidt, \emph{Diophantine approximation}, Lecture Notes in Mathematics 785, Springer-Verlag, Berlin Heidelberg, 1980.
\item[{[Sil]}] J. H. Silverman, \emph{The theory of height functions}, Arithmetic geometry (G. Cornell and J. H. Silverman, eds.), Springer-Verlag, New York, 1986, pp. 151--166.
\item[{[So]}] C. Soulé, \emph{Géométrie d'Arakelov et théorie des nombres transcendants}, Journées arithmétiques, Luminy, Astérisque（待刊）.
\item[{[So 2]}] C. Soulé, D. Abramovich, J.-F. Burnol, and J. Kramer, \emph{Lectures on Arakelov Geometry}, Cambridge studies in applied mathematics 33, Cambridge University Press, Cambridge, 1992.
\item[{[T]}] J. Tate, \emph{Rigid analytic spaces}, Invent. Math. 12 (1971), 257--289.
\item[{[U]}] K. Ueno, \emph{Classification theory of algebraic varieties and compact complex spaces}, Lecture Notes in Mathematics 439, Springer-Verlag, Berlin Heidelberg, 1975.
\item[{[V 1]}] P. Vojta, \emph{Diophantine approximations and value distribution theory}, Lecture Notes in Mathematics 1239, Springer-Verlag, Berlin Heidelberg, 1987.
\item[{[V 2]}] P. Vojta, \emph{Dyson's lemma for a product of two curves of arbitrary genus}, Invent. Math. 98 (1989), 107--113.
\item[{[V 3]}] P. Vojta, \emph{Mordell's conjecture over function fields}, Invent. Math. 98 (1989), 115--138.
\item[{[V 4]}] P. Vojta, \emph{Siegel's theorem in the compact case}, Ann. Math. 133 (1991), 509--548.
\item[{[V 5]}] P. Vojta, \emph{A generalization of theorems of Faltings and Thue-Siegel-Roth-Wirsing}, J. Amer. Math. Soc.（待刊）.
\end{description}
